\section{AMC 8 1988}
        \begin{enumerate}
        	\item (\textit{AMC 8 1988}) The diagram shows part of a scale of a measuring device.  The arrow indicates an approximate reading of $m$. What is the value of $m$?


 
                    \begin{figure}[H]
                        \centering
                        \includegraphics[width=0.5\textwidth]{Problem/AMC_8_AJHSME_1988_1_Problem_Figure_1.png}
                    \end{figure}
                    
$\text{(A)}\ 10.05 \qquad \text{(B)}\ 10.15 \qquad \text{(C)}\ 10.25 \qquad \text{(D)}\ 10.3 \qquad \text{(E)}\ 10.6$


	\item (\textit{AMC 8 1988}) The product $8\times .25\times 2\times .125 =$
$\text{(A)}\ \frac18 \qquad \text{(B)}\ \frac14 \qquad \text{(C)}\ \frac12 \qquad \text{(D)}\ 1 \qquad \text{(E)}\ 2$


	\item (\textit{AMC 8 1988}) Evaluate the following expression:


 $\frac{1}{10}+\frac{2}{20}+\frac{3}{30} =$
$\text{(A)}\ .1 \qquad \text{(B)}\ .123 \qquad \text{(C)}\ .2 \qquad \text{(D)}\ .3 \qquad \text{(E)}\ .6$


	\item (\textit{AMC 8 1988}) The figure consists of alternating light and dark squares.  The number of dark squares exceeds the number of light squares by


 $\text{(A)}\ 7 \qquad \text{(B)}\ 8 \qquad \text{(C)}\ 9 \qquad \text{(D)}\ 10 \qquad \text{(E)}\ 11$

                    \begin{figure}[H]
                        \centering
                        \includegraphics[width=0.5\textwidth]{Problem/AMC_8_AJHSME_1988_4_Problem_Figure_1.png}
                    \end{figure}
                    


	\item (\textit{AMC 8 1988}) If $\angle \text{CBD}$ is a right angle, then this protractor indicates that the measure of $\angle \text{ABC}$ is approximately


 
                    \begin{figure}[H]
                        \centering
                        \includegraphics[width=0.5\textwidth]{Problem/AMC_8_AJHSME_1988_5_Problem_Figure_1.png}
                    \end{figure}
                    
$\text{(A)}\ 20^\circ \qquad \text{(B)}\ 40^\circ \qquad \text{(C)}\ 50^\circ \qquad \text{(D)}\ 70^\circ \qquad \text{(E)}\ 120^\circ$


	\item (\textit{AMC 8 1988}) $\frac{(.2)^3}{(.02)^2} =$
$\text{(A)}\ .2 \qquad \text{(B)}\ 2 \qquad \text{(C)}\ 10 \qquad \text{(D)}\ 15 \qquad \text{(E)}\ 20$


	\item (\textit{AMC 8 1988}) $2.46\times 8.163\times (5.17+4.829)$ is closest to


 $\text{(A)}\ 100 \qquad \text{(B)}\ 200 \qquad \text{(C)}\ 300 \qquad \text{(D)}\ 400 \qquad \text{(E)}\ 500$


	\item (\textit{AMC 8 1988}) Betty used a calculator to find the product $0.075 \times 2.56$.  She forgot to enter the decimal points.  The calculator showed $19200$.  If Betty had entered the decimal points correctly, the answer would have been


 $\text{(A)}\ .0192 \qquad \text{(B)}\ .192 \qquad \text{(C)}\ 1.92 \qquad \text{(D)}\ 19.2 \qquad \text{(E)}\ 192$


	\item (\textit{AMC 8 1988}) An isosceles triangle is a triangle with two sides of equal length.  How many of the five triangles on the square grid below are isosceles?


 
                    \begin{figure}[H]
                        \centering
                        \includegraphics[width=0.5\textwidth]{Problem/AMC_8_AJHSME_1988_9_Problem_Figure_1.png}
                    \end{figure}
                    
$\text{(A)}\ 1 \qquad \text{(B)}\ 2 \qquad \text{(C)}\ 3 \qquad \text{(D)}\ 4 \qquad \text{(E)}\ 5$


	\item (\textit{AMC 8 1988}) Chris's birthday is on a Thursday this year.  What day of the week will it be $60$ days after her birthday?


 $\text{(A)}\ \text{Monday} \qquad \text{(B)}\ \text{Wednesday} \qquad \text{(C)}\ \text{Thursday} \qquad \text{(D)}\ \text{Friday} \qquad \text{(E)}\ \text{Saturday}$


	\item (\textit{AMC 8 1988}) $\sqrt{164}$ is 


 $\text{(A)}\ 42 \qquad \text{(B)}\ \text{less than }10 \qquad \text{(C)}\ \text{between }10\text{ and }11 \qquad \text{(D)}\ \text{between }11\text{ and }12 \qquad \text{(E)}\ \text{between }12\text{ and }13$


	\item (\textit{AMC 8 1988}) Suppose the estimated $20$ billion dollar cost to send a person to the planet Mars is shared equally by the $250$ million people in the U.S. Then each person's share is 


 $\text{(A)}\ 20\text{ dollars} \qquad \text{(B)}\ 50\text{ dollars} \qquad \text{(C)}\ 80\text{ dollars} \qquad \text{(D)}\ 100\text{ dollars} \qquad \text{(E)}\ 125\text{ dollars}$


	\item (\textit{AMC 8 1988}) If rose bushes are spaced about $1$ foot apart, approximately how many bushes are needed to surround a circular patio whose radius is $12$ feet?


 $\text{(A)}\ 12 \qquad \text{(B)}\ 38 \qquad \text{(C)}\ 48 \qquad \text{(D)}\ 75 \qquad \text{(E)}\ 450$


	\item (\textit{AMC 8 1988}) $\diamondsuit$ and $\Delta$ are whole numbers and $\diamondsuit \times \Delta =36$.  The largest possible value of $\diamondsuit + \Delta$ is 


 $\text{(A)}\ 12 \qquad \text{(B)}\ 13 \qquad \text{(C)}\ 15 \qquad \text{(D)}\ 20\ \qquad \text{(E)}\ 37$


	\item (\textit{AMC 8 1988}) The reciprocal of $\left( \frac{1}{2}+\frac{1}{3}\right)$ is 


 $\text{(A)}\ \frac{1}{6} \qquad \text{(B)}\ \frac{2}{5} \qquad \text{(C)}\ \frac{6}{5} \qquad \text{(D)}\ \frac{5}{2} \qquad \text{(E)}\ 5$


	\item (\textit{AMC 8 1988}) Placing no more than one $\text{X}$ in each small square, what is the greatest number of $\text{X}$'s that can be put on the grid shown without getting three $\text{X}$'s in a row vertically, horizontally, or diagonally?


 $\text{(A)}\ 2 \qquad \text{(B)}\ 3 \qquad \text{(C)}\ 4 \qquad \text{(D)}\ 5 \qquad \text{(E)}\ 6$

                    \begin{figure}[H]
                        \centering
                        
\includegraphics[width=0.5\textwidth]{Problem/AMC_8_AJHSME_1988_16_Problem_Figure_1.png}
                    \end{figure}
                    


	\item (\textit{AMC 8 1988}) The shaded region formed by the two intersecting perpendicular rectangles, in square units, is


 
                    \begin{figure}[H]
                        \centering
                        \includegraphics[width=0.5\textwidth]{Problem/AMC_8_AJHSME_1988_17_Problem_Figure_1.png}
                    \end{figure}
                    
$\text{(A)}\ 23 \qquad \text{(B)}\ 38 \qquad \text{(C)}\ 44 \qquad \text{(D)}\ 46 \qquad \text{(E)}\ \text{unable to be determined from the information given}$


	\item (\textit{AMC 8 1988}) The average weight of $6$ boys is $150$ pounds and the average weight of $4$ girls is $120$ pounds.  The average weight of the $10$ children is 


 $\text{(A)}\ 135\text{ pounds} \qquad \text{(B)}\ 137\text{ pounds} \qquad \text{(C)}\ 138\text{ pounds} \qquad \text{(D)}\ 140\text{ pounds} \qquad \text{(E)}\ 141\text{ pounds}$


	\item (\textit{AMC 8 1988}) What is the $100\text{th}$ number in the arithmetic sequence: $1,5,9,13,17,21,25,...$?


 $\text{(A)}\ 397 \qquad \text{(B)}\ 399 \qquad \text{(C)}\ 401 \qquad \text{(D)}\ 403 \qquad \text{(E)}\ 405$


	\item (\textit{AMC 8 1988}) The glass gauge on a cylindrical coffee maker shows that there are $45$ cups left when the coffee maker is $36\%$ full.  How many cups of coffee does it hold when it is full?


 $\text{(A)}\ 80 \qquad \text{(B)}\ 100 \qquad \text{(C)}\ 125 \qquad \text{(D)}\ 130 \qquad \text{(E)}\ 262$

                    \begin{figure}[H]
                        \centering
                        \includegraphics[width=0.5\textwidth]{Problem/AMC_8_AJHSME_1988_20_Problem_Figure_1.png}
                    \end{figure}
                    


	\item (\textit{AMC 8 1988}) A fifth number, $n$, is added to the set $\{ 3,6,9,10 \}$ to make the mean of the set of five numbers equal to its median.  The number of possible values of $n$ is


 $\text{(A)}\ 1 \qquad \text{(B)}\ 2 \qquad \text{(C)}\ 3 \qquad \text{(D)}\ 4 \qquad \text{(E)}\ \text{more than }4$


	\item (\textit{AMC 8 1988}) Tom's Hat Shoppe increased all original prices by $25\%$.  Now the shoppe is having a sale where all prices are $20\%$ off these increased prices.  Which statement best describes the sale price of an item?


 $\text{(A)}\ \text{The sale price is }5\% \text{ higher than the original price.}$
$\text{(B)}\ \text{The sale price is higher than the original price, but by less than }5\% .$
$\text{(C)}\ \text{The sale price is higher than the original price, but by more than }5\% .$
$\text{(D)}\ \text{The sale price is lower than the original price.}$
$\text{(E)}\ \text{The sale price is the same as the original price.}$


	\item (\textit{AMC 8 1988}) Maria buys computer disks at a price of $4$ for $<dollar>5$ and sells them at a price of $3$ for $<dollar>$$5$.  How many computer disks must she sell in order to make a profit of $<dollar>100$?


 $\text{(A)}\ 100 \qquad \text{(B)}\ 120 \qquad \text{(C)}\ 200 \qquad \text{(D)}\ 240 \qquad \text{(E)}\ 1200$


	\item (\textit{AMC 8 1988}) 
                    \begin{figure}[H]
                        \centering
                        \includegraphics[width=0.5\textwidth]{Problem/AMC_8_AJHSME_1988_24_Problem_Figure_1.png}
                    \end{figure}
                    


 The square in the first diagram "rolls" clockwise around the fixed regular hexagon until it reaches the bottom.  In which position will the solid triangle be in diagram $4$?


 
                    \begin{figure}[H]
                        \centering
                        
\includegraphics[width=0.5\textwidth]{Problem/AMC_8_AJHSME_1988_24_Problem_Figure_2.png}
                    \end{figure}
                    


	\item (\textit{AMC 8 1988}) A \textbf{palindrome} is a whole number that reads the same forwards and backwards.  If one neglects the colon, certain times displayed on a digital watch are palindromes.  Three examples are: $\boxed{1:01}$, $\boxed{4:44}$, and $\boxed{12:21}$.  How many times during a $12$-hour period will be palindromes?


 $\text{(A)}\ 57 \qquad \text{(B)}\ 60 \qquad \text{(C)}\ 63 \qquad \text{(D)}\ 90 \qquad \text{(E)}\ 93$


\end{enumerate}